% ============================================================================
% CHƯƠNG 4: THIẾT LẬP THỰC NGHIỆM
% ============================================================================
\chapter{Thiết lập thực nghiệm}
\label{chap:experiment_setup}

Chương này mô tả chi tiết môi trường thực nghiệm, tập dữ liệu, cấu hình huấn luyện và các mô hình được sử dụng để so sánh với \fastvit{}.

% ============================================================================
\section{Môi trường thực nghiệm}
\label{sec:environment}

% ----------------------------------------------------------------------------
\subsection{Phần cứng}
\label{subsec:hardware}

\begin{table}[H]
\centering
\caption{Cấu hình phần cứng thực nghiệm.}
\label{tab:hardware}
\begin{tabular}{l l}
\toprule
\textbf{Thành phần} & \textbf{Chi tiết} \\
\midrule
GPU & NVIDIA H100 80GB HBM3 \\
Số CPU & 20 (logical: 20) \\
RAM hệ thống & $\geq$ 32 GB \\
Lưu trữ & SSD NVMe \\
Mobile Benchmark (gốc) & iPhone 12 Pro (Apple A14 Bionic) \\
\bottomrule
\end{tabular}
\end{table}

% ----------------------------------------------------------------------------
\subsection{Phần mềm}
\label{subsec:software}

\begin{table}[H]
\centering
\caption{Cấu hình phần mềm và thư viện.}
\label{tab:software}
\begin{tabular}{l l l}
\toprule
\textbf{Thành phần} & \textbf{Phiên bản} & \textbf{Vai trò} \\
\midrule
Python & 3.9 & Ngôn ngữ lập trình \\
PyTorch & 1.11.0 & Framework deep learning \\
TorchVision & 0.12.0 & Tiện ích xử lý ảnh \\
timm & $\geq$ 0.6.x & Thư viện mô hình, data augmentation \\
fvcore & -- & Đếm FLOPs và tham số \\
cuDNN & auto & Tối ưu convolution trên GPU \\
CUDA Toolkit & 11.3 & Hỗ trợ tính toán GPU \\
Wandb & -- & Logging và monitoring \\
\bottomrule
\end{tabular}
\end{table}

% ============================================================================
\section{Tập dữ liệu}
\label{sec:dataset}

Thực nghiệm phát hiện đối tượng (Object Detection) sử dụng tập dữ liệu \textbf{MS-COCO 2017}:

\begin{table}[H]
\centering
\caption{Thông tin tập dữ liệu MS-COCO 2017.}
\label{tab:dataset_coco}
\begin{tabular}{l r}
\toprule
\textbf{Thuộc tính} & \textbf{Giá trị} \\
\midrule
Số lớp (categories) & 80 \\
Ảnh huấn luyện (train2017) & $\sim$118{,}287 \\
Ảnh kiểm định (val2017) & 5{,}000 \\
Số annotations (train) & $\sim$860{,}000 instances \\
Kích thước ảnh & Đa dạng (resize về $800 \times 800$ khi huấn luyện) \\
Định dạng annotation & JSON (COCO format) \\
\bottomrule
\end{tabular}
\end{table}

MS-COCO là benchmark tiêu chuẩn cho phát hiện đối tượng, bao gồm 80 lớp đối tượng đa dạng từ ``person'', ``car'', ``cat'' đến các vật dụng nhỏ như ``toothbrush'', ``scissors''. So với ImageNet (chỉ phân loại), COCO đòi hỏi mô hình phải định vị chính xác (bounding box regression) ngoài phân loại.

% ----------------------------------------------------------------------------
\subsection{Phân tích đặc trưng dữ liệu (EDA)}
\label{subsec:coco_eda}

Tập dữ liệu MS-COCO có những đặc trưng phân phối đặc biệt, tạo ra những thách thức lớn đối với các kiến trúc Backbone và Detector. So với các bộ dữ liệu thế hệ trước như PASCAL VOC, MS-COCO tập trung vào sự hiểu biết ngữ cảnh phức tạp của các vật thể trong môi trường tự nhiên.

\textbf{1. Phân phối kích thước đối tượng (Object Scale Distribution)}

MS-COCO phân loại các đối tượng thành 3 nhóm dựa trên diện tích pixel. Điểm đặc trưng nhất của COCO là sự thống trị của các vật thể kích thước siêu nhỏ.

\begin{figure}[H]
\centering
\begin{tikzpicture}
\begin{axis}[
    ybar,
    ylabel={Tỷ lệ phần trăm (\%)},
    symbolic x coords={Nhỏ ($<32^2$), Trung bình ($32^2-96^2$), Lớn ($>96^2$)},
    xtick=data,
    ymin=0, ymax=50,
    bar width=25pt,
    width=0.6\textwidth,
    height=6cm,
    nodes near coords,
    every node near coord/.append style={font=\small},
]
\addplot[fill=blue!40] coordinates {
    (Nhỏ ($<32^2$), 41) 
    (Trung bình ($32^2-96^2$), 34) 
    (Lớn ($>96^2$), 24)
};
\end{axis}
\end{tikzpicture}
\caption{Phân phối kích thước đối tượng trong MS-COCO.}
\label{fig:coco_scale}
\end{figure}

Sự áp đảo của các vật thể nhỏ (hơn 41\%) giải thích vì sao Mask R-CNN đòi hỏi ảnh đầu vào phải được resize lên kích thước lớn ($800 \times 800$) thay vì $256 \times 256$ như bài toán phân loại. Khi qua các lớp downsampling của backbone, các vật thể này rất dễ bị ``xóa sổ'' trên feature map, do đó kiến trúc Feature Pyramid Network (FPN) là bắt buộc để khôi phục đặc trưng ở nhiều mức độ phân giải khác nhau.

\textbf{2. Mật độ vật thể và hiện tượng che khuất (Instance Density \& Occlusion)}

Khác với các bộ dữ liệu mà vật thể thường nằm ở trung tâm và có background rõ ràng, MS-COCO chứa các bối cảnh đời thực với mật độ rất dày đặc:
\begin{itemize}
    \item Trung bình có \textbf{7.3 đối tượng trên mỗi bức ảnh}.
    \item Hơn 10\% số ảnh có trên 20 đối tượng.
    \item Hiện tượng che khuất (occlusion) xảy ra thường xuyên khi các đối tượng chồng chéo lên nhau (ví dụ: người đứng trong đám đông, các phương tiện nối đuôi nhau). Điều này đặt ra yêu cầu rất cao đối với độ phân giải không gian của feature map và khả năng trích xuất ranh giới của mạng.
\end{itemize}

\textbf{3. Mất cân bằng lớp (Long-Tailed Class Imbalance)}

Phân phối nhãn trong COCO cực kỳ bất đối xứng, tuân theo phân phối đuôi dài (long-tailed distribution). Bảng dưới đây minh họa sự chênh lệch giữa các lớp xuất hiện nhiều nhất (Head) và ít nhất (Tail):

\begin{table}[H]
\centering
\caption{Sự mất cân bằng giữa các lớp (Head vs Tail) trong MS-COCO.}
\label{tab:coco_imbalance}
\begin{tabular}{l r | l r}
\toprule
\multicolumn{2}{c|}{\textbf{Top 5 lớp nhiều nhất (Head)}} & \multicolumn{2}{c}{\textbf{Top 5 lớp ít nhất (Tail)}} \\
\textbf{Lớp (Class)} & \textbf{Số lượng (xấp xỉ)} & \textbf{Lớp (Class)} & \textbf{Số lượng (xấp xỉ)} \\
\midrule
person & $\sim$262{,}000 & hair drier & $\sim$190 \\
car & $\sim$44{,}000 & toaster & $\sim$225 \\
chair & $\sim$38{,}000 & parking meter & $\sim$700 \\
book & $\sim$24{,}000 & bear & $\sim$1{,}300 \\
bottle & $\sim$24{,}000 & scissors & $\sim$1{,}500 \\
\bottomrule
\end{tabular}
\end{table}

Sự chênh lệch khổng lồ (lớp \texttt{person} nhiều gấp $\sim 1{,}300$ lần lớp \texttt{hair drier}) gây ra hiện tượng thiên lệch (bias) trong quá trình học. Các lớp Head dễ dàng đạt Average Precision (AP) cao do có lượng dữ liệu dồi dào, trong khi các lớp Tail lại cực kỳ thấp do mô hình không có đủ các mẫu đa dạng để học được các bất biến (invariants).

\textbf{4. Độ phức tạp của ngữ cảnh (Context Complexity)}

Các vật thể trong COCO hiếm khi xuất hiện đơn độc mà luôn gắn liền với ngữ cảnh tự nhiên (người cưỡi ngựa, cái nĩa nằm cạnh đĩa thức ăn, cái cốc trên bàn). Sự tương tác giữa các vật thể đòi hỏi backbone phải có cơ chế tiếp nhận ngữ cảnh toàn cục (\textit{global context}) mạnh mẽ để suy luận mối quan hệ giữa chúng, từ đó giảm thiểu dự đoán sai (False Positives). Đây cũng là lý do tại sao các kiến trúc Hybrid kết hợp Convolution (nắm bắt chi tiết cục bộ) và Attention (hiểu ngữ cảnh toàn cục) như \fastvit{} lại phát huy được thế mạnh rất lớn trên tập dữ liệu này.

% ============================================================================
\section{Tiền xử lý dữ liệu}
\label{sec:preprocessing}

Quy trình tiền xử lý dữ liệu tuân theo cấu hình tiêu chuẩn cho Mask R-CNN trên COCO:

\textbf{Huấn luyện (Training):}
\begin{itemize}
    \item \textbf{Resize:} Thay đổi kích thước ảnh sao cho cạnh ngắn đạt $800$ pixel, duy trì aspect ratio.
    \item \textbf{Random Horizontal Flip:} Xác suất 50\% để tăng cường dữ liệu.
    \item \textbf{Chuẩn hóa:} Theo ImageNet mean $= (0.485, 0.456, 0.406)$ và std $= (0.229, 0.224, 0.225)$ do backbone được pre-train trên ImageNet.
\end{itemize}

\textbf{Kiểm định (Validation):}
\begin{itemize}
    \item \textbf{Resize:} Kích thước ảnh $800 \times 800$.
    \item \textbf{Chuẩn hóa:} Giống huấn luyện.
\end{itemize}

% ============================================================================
\section{Cấu hình huấn luyện}
\label{sec:training_config}

Thực nghiệm sử dụng kiến trúc \textbf{Mask R-CNN} với backbone \fastvit{}-SA12 được pre-train, huấn luyện trên MS-COCO 2017.

\begin{table}[H]
\centering
\caption{Cấu hình huấn luyện Mask R-CNN + \fastvit{}-SA12 trên MS-COCO.}
\label{tab:training_config_det}
\begin{tabular}{l l}
\toprule
\textbf{Siêu tham số} & \textbf{Giá trị} \\
\midrule
Kiến trúc detector & Mask R-CNN \\
Backbone & FastViT-SA12 (pre-trained ImageNet-1K) \\
FPN channels & 256 \\
Optimizer & AdamW \\
Learning rate (LR) & $10^{-4}$ \\
Weight decay & 0.05 \\
LR scheduler & Step (giảm tại epoch 8 và 11) \\
LR gamma & 0.1 \\
Warmup & 5 epochs \\
Tổng epochs & 12 (schedule 1$\times$) \\
Batch size & 16 \\
Gradient accumulation & 1 step \\
Gradient clipping & 5.0 \\
Input size & $800 \times 800$ \\
Mixed Precision & Native AMP (FP16) \\
RPN pre-NMS (train) & 1{,}000 proposals \\
RPN post-NMS (train) & 1{,}000 proposals \\
RPN batch size & 128 \\
Box batch size & 128 \\
Workers & 15 \\
Phần cứng & NVIDIA H100 80GB HBM3, 20 CPUs \\
\bottomrule
\end{tabular}
\end{table}

\textbf{Phân tích cấu hình detection:}
\begin{itemize}
    \item \textbf{Mask R-CNN} là kiến trúc two-stage detector tiêu chuẩn, bao gồm: Region Proposal Network (RPN) $\to$ ROI Align $\to$ Classification/Regression heads. Đây là framework đánh giá backbone phổ biến trong cộng đồng.
    
    \item \textbf{Learning rate $10^{-4}$} do backbone đã pre-train, chỉ cần fine-tune.
    
    \item \textbf{Step scheduler} (giảm LR $\times 0.1$ tại epoch 8 và 11) theo schedule ``1$\times$'' tiêu chuẩn cho COCO detection.
    
    \item \textbf{FPN 256 channels:} Feature Pyramid Network kết hợp feature maps đa tỷ lệ từ backbone \fastvit{}-SA12 (4 stages), cho phép phát hiện đối tượng ở nhiều kích thước.
    
    \item \textbf{Input $800 \times 800$:} Kích thước cần thiết cho detection vì các đối tượng nhỏ đòi hỏi độ phân giải cao để trích xuất đặc trưng chính xác.
    
    \item \textbf{Gradient clipping 5.0:} Ổn định huấn luyện khi sử dụng AMP trên mô hình phức tạp (backbone + FPN + heads).
\end{itemize}

\textbf{Command huấn luyện:}
\begin{lstlisting}[language=bash, basicstyle=\ttfamily\scriptsize, breaklines=true]
python object_detection.py --arch maskrcnn \
  --model fastvit_sa12 \
  --pretrained-backbone fastvit_sa12.pth.tar \
  --dataset coco --data-dir ./data/coco \
  --epochs 12 --batch-size 16 --lr 1e-4 \
  --lr-steps 8 11 --lr-gamma 0.1 \
  --weight-decay 0.05 --clip-grad 5.0 \
  --fpn-channels 256 --amp --scheduler step
\end{lstlisting}

% ============================================================================
\section{Mô hình đối chiếu}
\label{sec:comparison_models}

Do tập trung vào thực nghiệm MS-COCO với biến thể \textbf{FastViT-SA12}, báo cáo sẽ sử dụng \textbf{ResNet-50} (làm backbone cho Mask R-CNN) như là mô hình tham chiếu cơ sở để so sánh. ResNet-50 là backbone phổ biến nhất cho Mask R-CNN, cho phép đánh giá hiệu quả vượt trội của FastViT-SA12.

\begin{itemize}
    \item \textbf{FastViT-SA12:} $\sim$11M tham số.
    \item \textbf{ResNet-50:} $\sim$25.6M tham số.
\end{itemize}

% ============================================================================
\section{Các chỉ số đánh giá}
\label{sec:metrics}

\textbf{Chỉ số đánh giá phát hiện đối tượng:}

\begin{description}
    \item[AP (Average Precision):] Diện tích dưới đường Precision-Recall cho mỗi lớp đối tượng tại một ngưỡng IoU cụ thể.
    
    \item[mAP (mean Average Precision):] Trung bình AP trên tất cả các lớp:
    \begin{equation}
        \text{mAP} = \frac{1}{C} \sum_{c=1}^{C} \text{AP}_c
    \end{equation}
    
    \item[mAP@0.50:] mAP tại ngưỡng IoU $= 0.50$ (tiêu chuẩn PASCAL VOC).
    
    \item[mAP@{[}0.5:0.95{]}:] Trung bình mAP tại 10 ngưỡng IoU từ 0.50 đến 0.95 (bước 0.05). Đây là chỉ số chính thức của COCO, khắt khe hơn mAP@0.50 vì yêu cầu bounding box chính xác hơn:
    \begin{equation}
        \text{mAP@[0.5:0.95]} = \frac{1}{10} \sum_{t \in \{0.50, 0.55, \ldots, 0.95\}} \text{mAP@}t
    \end{equation}
    
    \item[IoU (Intersection over Union):] Độ trùng lắp giữa bounding box dự đoán và ground truth:
    \begin{equation}
        \text{IoU} = \frac{|B_{\text{pred}} \cap B_{\text{gt}}|}{|B_{\text{pred}} \cup B_{\text{gt}}|}
    \end{equation}
\end{description}

\textbf{Chỉ số đánh giá hiệu quả:}

\begin{description}
    \item[Parameters (Params):] Tổng số tham số có thể học. Đơn vị: triệu (M).
    
    \item[FLOPs (Floating Point Operations):] Số phép nhân--cộng dấu phẩy động cho một lần inference. Đơn vị: GigaFLOPs (G). Đo bằng thư viện \texttt{fvcore}. FLOPs phản ánh chi phí tính toán lý thuyết, không phụ thuộc phần cứng.
    
    \item[Latency:] Thời gian thực tế cho một lần suy luận. Đơn vị: milliseconds (ms). Đo trên phần cứng cụ thể (GPU hoặc Mobile). Latency phản ánh hiệu quả thực tế, bao gồm cả overhead phần cứng.
    
    \item[Throughput:] Số ảnh xử lý mỗi giây. Đơn vị: images/second (img/s).
    \begin{equation}
        \text{Throughput} = \frac{\text{Batch size}}{\text{Latency (s)}}
    \end{equation}
\end{description}

\textbf{Mối quan hệ FLOPs vs. Latency:}

FLOPs và Latency không luôn tỷ lệ thuận. Các yếu tố ảnh hưởng latency ngoài FLOPs:
\begin{itemize}
    \item \textbf{Memory access cost:} Các phép toán memory-bound (reshape, concat, softmax) có latency cao hơn FLOPs gợi ý.
    \item \textbf{Mức độ song song hóa:} Depthwise conv có FLOPs thấp nhưng arithmetic intensity thấp, chưa chắc nhanh hơn dense conv trên GPU.
    \item \textbf{Số lượng operators:} Nhiều operators nhỏ liên tiếp có overhead kernel launch trên GPU.
    \item \textbf{Reparameterization:} Giảm số operators, giảm memory access, cải thiện latency vượt xa mức giảm FLOPs.
\end{itemize}

Đây là lý do \fastvit{} tập trung tối ưu \textit{latency thực tế} (qua reparameterization) thay vì chỉ giảm FLOPs.
